\usepackage[magyar]{babel}
\usepackage{fontspec}
\setmonofont{DejaVu Sans Mono}[Scale=0.80]

% --- Színek (fehér alap, fekete szöveg, színes kiemelések) ---
\usepackage{xcolor}
\definecolor{linkblue}{RGB}{0,90,180}
\definecolor{headcolor}{RGB}{30,80,140}
\definecolor{keybox}{RGB}{240,245,250}
\definecolor{keyborder}{RGB}{0,90,180}
\definecolor{codebg}{RGB}{245,245,245}

% --- Linkek színezése ---
\usepackage{hyperref}
\hypersetup{
  colorlinks=true,
  linkcolor=linkblue,
  urlcolor=linkblue,
  citecolor=linkblue,
  filecolor=linkblue,
}

% --- Fejléc/lábléc ---
\usepackage{fancyhdr}
\pagestyle{fancy}
\fancyhead[L]{\small\color{gray} Az Ágensek Kora}
\fancyhead[R]{\small\color{gray}\thepage}
\fancyfoot[C]{}
\renewcommand{\headrulewidth}{0.4pt}
\renewcommand{\headrule}{\hbox to\headwidth{\color{gray!50}\leaders\hrule height \headrulewidth\hfill}}

% --- Fejezet/szekció színek ---
\usepackage{titlesec}
\titleformat{\chapter}[display]
  {\normalfont\huge\bfseries\color{headcolor}}
  {\chaptertitlename\ \thechapter}{20pt}{\Huge}
\titleformat{\section}
  {\normalfont\Large\bfseries\color{headcolor}}
  {\thesection}{1em}{}
\titleformat{\subsection}
  {\normalfont\large\bfseries\color[RGB]{60,60,60}}
  {\thesubsection}{1em}{}

% --- Kulcs üzenet box ---
\usepackage{tcolorbox}
\newtcolorbox{kulcsuzenat}{
  colback=keybox,
  colframe=keyborder,
  coltext=black,
  title={\textcolor{white}{Kulcs üzenet PM-nek}},
  coltitle=white,
  colbacktitle=keyborder,
  fonttitle=\bfseries,
  boxrule=1pt,
  arc=4pt,
  left=8pt,
  right=8pt,
  top=6pt,
  bottom=6pt,
}

% --- Táblázatok: kisebb betűméret ---
\usepackage{longtable}
\usepackage{booktabs}
\usepackage{etoolbox}
\AtBeginEnvironment{longtable}{\small}
\setlength{\LTpre}{0.5em plus 5pt minus 3pt}
\setlength{\LTpost}{0.5em plus 3pt minus 2pt}
% Szűkebb sorok a táblázatokban
\renewcommand{\arraystretch}{0.9}

% --- Kód blokk háttér ---
\usepackage{fancyvrb}
\usepackage{framed}
\definecolor{shadecolor}{RGB}{245,245,245}

% --- Magyar nevek ---
\renewcommand{\contentsname}{Tartalomjegyzék}
\renewcommand{\chaptername}{Fejezet}

% --- Tartalomjegyzék ---
\usepackage{tocloft}
\renewcommand{\cfttoctitlefont}{\huge\bfseries\color{headcolor}}
\renewcommand{\cftchapfont}{\bfseries\color{black}}
\renewcommand{\cftsecfont}{\color[RGB]{60,60,60}}
\renewcommand{\cftchappagefont}{\color{black}}
\renewcommand{\cftsecpagefont}{\color[RGB]{60,60,60}}
